\section{RQ2--RQ4 Findings and Evidence Status}

For RQ2, the 90-day no-repeat-purchase target was reconstructed from order
headers with complete agreement. The selected CatBoost classifier achieved
TEST PR-AUC 0.7461, ROC-AUC 0.8314 and Brier score 0.1650. A revised
training-only operating threshold increased precision to 0.6064, retained
recall of 0.8866 and raised F1-score to 0.7202. False positives fell from
7,169 to 5,016 (30.0\%), while false negatives rose from 350 to 988. This
is a threshold trade-off rather than a material gain in ranking ability.
The TEST period had already been examined in the preceding experiment, so
the revised result is a repeated-holdout diagnostic rather than independent
confirmation. These results support temporal risk ranking for the existing
customer population; they do not establish permanent churn, cold-start
performance or the causal effect of contacting customers.

For RQ3, the supplied trend-policy scenario had lower scenario cost and
average stock but higher lost units than the baseline. Because future demand
is synthetic oracle demand, the cost formula is undocumented and the participating
business's observed replenishment policy was not supplied, the result is a scenario trade-off rather
than evidence of realised savings or policy improvement.

For RQ4, the encoded replenishment rule reproduced the supplied quantity,
action and constraint reason in all prepared cases. The new English and
Bangla explanation generator included the required numerical evidence in all
cases. The constant \texttt{human\_rating} field is not participant feedback;
therefore the research question concerning task accuracy, confidence and
perceived usability among SME owners and managers cannot yet be answered.
A counterbalanced study involving SME owners and managers is required.

The next validation stage is an untouched future and held-out-business
evaluation for predictive models, a documented inventory replay against an
observed replenishment policy of a participating business, and a consent-based
study involving SME owners and managers to evaluate the explanations.
Until these validation requirements are met, the models and recommendations
should be treated as research outputs for decision support, not as autonomous
business actions.
